\begin{textsample}{NEG Dim 5 – global\_north\_2024\_03 – Score 1.00 – global\_north\_2024\_03/106325716.txt}  \label{ex:f5_neg_090}
Palestinian factions including rivals Hamas and Fatah said on Friday they would pursue"unity of action"in confronting Israel after representatives met at Russia-hosted talks.
The meeting in Moscow on Thursday brought together Hamas, Islamic Jihad, Fatah and other Palestinian groups for talks on the war in Gaza and an eventual post-war period.
It came on the heels of the resignation of the Palestinian Authority government, which is led by Fatah and based in the occupied West Bank.
Outgoing prime minister Mohammad Shtayyeh called for intra-Palestinian consensus as he announced the resignation, and some analysts said the development could pave the way for a government of technocrats that could operate in the West Bank and Hamas-run Gaza after the war.
Arab and Western leaders have been pushing for reforms to the Palestinian Authority as they discuss possible reconstruction efforts.
A statement on Friday by the Palestinian factions represented in Moscow said there would be an"upcoming dialogue"to bring them under the banner of the Palestine Liberation Organization ( PLO ).
Thursday ' s"constructive"talks saw @ @ @ @ @ @ @ @ @ @ Israeli forces from Gaza and the creation of a Palestinian state, the statement said.
While Hamas and Islamic Jihad are considered"terrorist"entities by Western powers, the PLO is internationally recognised as representing Palestinians in the Palestinian territories and diaspora.
Discussions in recent years about integrating Hamas into the PLO have ended in failure.
In recent years, Moscow has strived to maintain good relations with all actors in the Israeli-Palestinian conflict, including Fatah and Hamas.
Russia ' s relations with Israel have become strained amid Moscow ' s criticism of Israeli actions in Gaza and rejection of a Palestinian state.
The war in Gaza was triggered by Hamas ' s October 7 attack on Israel that resulted in the deaths of about 1,160 people, mostly civilians, according to an AFP tally based on official Israeli figures.
At least 30,228 people, mostly women and children, have been killed in Israel ' s retaliatory military offensive in Gaza, according to the health ministry in the Hamas-run territory.
Disclaimer : The content @ @ @ @ @ @ @ @ @ @ from an external third party provider.
We are not responsible for, and do not control, such external websites, entities, applications or media publishers.
The body of the text is provided on an"as is"and"as available"basis and has not been edited in any way.
Neither we nor our affiliates guarantee the accuracy of or endorse the views or opinions expressed in this article.
Read our full disclaimer policy here.

% matched lemmas: 
\end{textsample}
